\documentclass[tikz,border=6pt]{standalone}
\usepackage{peabra}
\begin{document}
\begin{tikzpicture}[node distance=4mm and 6mm, every node/.append style={font=\footnotesize}]
  \node[sotetdoboz=34mm, font=\footnotesize\bfseries] (a) {1. állapotok\\definiálása};
  \node[doboz=34mm, right=of a] (b) {2. átmenetek és ráták ($\lambda$, $\mu$, $\beta\lambda$, \dots)};
  \node[doboz=34mm, right=of b] (c) {3. a $\mathbf{Q}$ intenzitás\-mátrix felírása};
  \node[doboz=30mm, below=10mm of a] (d1) {tranziens: $\dot{\mathbf{p}} = \mathbf{p}\mathbf{Q}$\\$\mathbf{p}(t) = \mathbf{p}(0)\,e^{\mathbf{Q}t}$};
  \node[doboz=30mm, right=of d1] (d2) {állandósult: $\boldsymbol{\pi}\mathbf{Q} = 0$, $\sum \pi_i = 1$};
  \node[doboz=30mm, right=of d2] (d3) {elnyelő állapot: MTTF az alapmátrixból};
  \node[doboz=30mm, right=of d3] (d4) {többfázisú: a próbáknál az állapotok újrarendezése};
  \node[zolddoboz=60mm, below=10mm of d2, xshift=17mm] (e) {4. eredmények: $A(t)$, $A_\infty$, MTTF, PFD$(t)$, PFD$_{\text{avg}}$, PFH};
  \node[szurkedoboz=60mm, below=of e] (f) {5. érzékenység és bizonytalanság};
  \draw[nyil] (a) -- (b); \draw[nyil] (b) -- (c);
  \draw[nyil] (c.south) -- ++(0,-0.35) -| (d1.north); \draw[nyil] (c.south) -- ++(0,-0.35) -| (d2.north);
  \draw[nyil] (c.south) -- ++(0,-0.35) -| (d3.north); \draw[nyil] (c.south) -- ++(0,-0.35) -| (d4.north);
  \foreach \d in {d1,d2,d3,d4} \draw[nyil] (\d.south) -- ++(0,-0.3) -| (e.north);
  \draw[nyil] (e) -- (f);
\end{tikzpicture}
\end{document}
